\section{O desenho da política e a sua teoria da mudança}
\label{sec:desenho}

Avaliar o desenho de uma política é examinar a lógica que liga o problema, a
ação e o resultado esperado, e explicitar as hipóteses que cada passo dessa
lógica pressupõe \citep{gertler2016}. A lei cearense não traz essa cadeia
escrita; ela se reconstrói da justificativa do projeto de lei e do debate que o
cercou.

\subsection{A hipótese causal}

\emph{Se} a lei proíbe a aplicação por aeronave e a proibição é cumprida,
\emph{então} cessa a deriva de fungicidas sobre comunidades e mananciais, o que
reduz a exposição de gestantes e, ao final, melhora os desfechos de nascimento,
contribuindo para reduzir a externalidade da fruticultura irrigada sobre a
população do entorno.

\begin{table}[ht]
\centering\small
\caption{Cadeia de resultados da proibição, com indicadores, premissas (P) e
riscos (R)}
\label{tab:tdm}
\begin{tabularx}{\textwidth}{@{}p{2.3cm}LLL@{}}
\toprule
Elo & Conteúdo & Indicador e fonte & Premissas e riscos \\
\midrule
Problema & Deriva de fungicidas aplicados por avião alcança casas, escolas e água & Resíduos em água; relatos de campo & P: a prática existia às vésperas da lei \\
Atividade & Vedação da aplicação aérea, com multa & Texto legal; autos de infração & P: há fiscalização \\
Produto & Deixam de ocorrer voos de aplicação & Relatórios mensais ao Ministério da Agricultura & R: descumprimento \\
Resultado intermediário & Menor exposição ambiental de gestantes & Não observado diretamente & R: substituição por aplicação terrestre, com mais calda e mais aplicadores \\
Resultado final, saúde & Maior peso ao nascer; menos prematuridade e perda fetal & SINASC, SIM & R: seleção para nascimento vivo \\
Resultado final, economia & Custo de controle da doença; área e produção & PAM, RAIS & R: migração da produção para o estado vizinho \\
\bottomrule
\end{tabularx}
\end{table}

\subsection{Avaliação do desenho}

\paragraph{O problema está bem definido, mas o diagnóstico é anterior.} A
carência que a lei enfrenta, a exposição involuntária de vizinhos, tem
população-alvo identificável e base empírica. Essa base, porém, descreve a
década de 2000. O legislador de 2018 decidiu com o diagnóstico de 2010, e a
primeira premissa da cadeia, a de que a prática estava em curso, não foi
verificada no processo legislativo. É a premissa que este projeto mede
(Seção~\ref{sec:quando}).

\paragraph{O instrumento ataca o método, não a molécula.} A lei proíbe uma forma
de aplicar. Os mesmos produtos seguem permitidos por via terrestre. A cadeia só
se completa se a deriva da aplicação terrestre for menor que a do avião, o que
é plausível, e se a troca não criar outra exposição. Produtores relataram à
imprensa que o volume de calda passou de 25 para 250 a 300 litros por hectare e
que a aplicação manual ocupa dezenas de trabalhadores onde havia um piloto. O
risco é que a exposição ambiental caia e a ocupacional suba. A teoria da mudança
gera, assim, duas previsões de sinais opostos, sobre nascimentos e sobre
intoxicações de trabalhadores, e a avaliação testa as duas.

\paragraph{O funil de atrito é estreito.} De cada cem nascimentos nos municípios
atingidos, poucos são de mães que moravam ao alcance da deriva. O efeito final
sobre a média municipal é o efeito sobre as expostas multiplicado por essa
fração. Superestimar a participação leva a amostras pequenas demais para
detectar o efeito \citep{gertler2016}; a Seção~\ref{sec:poder} faz a conta.

\paragraph{A fiscalização ficou mais simples.} A regra anterior exigia medir
distâncias e condições de voo. A proibição total dispensa a medição. É o
argumento mais forte a favor do desenho da lei, e independe do tamanho do
efeito sobre a saúde.

\paragraph{Não há plano de monitoramento.} A lei não previu indicadores, linha
de base nem avaliação. Os dados que permitiriam acompanhá-la existem em três
órgãos e nunca foram cruzados.

\subsection{Seis caminhos entre a lei e o nascimento}
\label{sec:caminhos}

A Figura~\ref{fig:dag} representa a cadeia como grafo causal. O que uma
avaliação de impacto estima é o efeito \emph{total} da lei sobre o desfecho, por
todos os caminhos. Só um deles é a hipótese da política.

\begin{figure}[ht]
\centering
\begin{tikzpicture}[>=Stealth, node distance=7mm and 13mm,
  every node/.style={font=\footnotesize},
  obs/.style={draw, rounded corners=2pt, align=center, inner sep=3pt},
  lat/.style={draw, dashed, rounded corners=2pt, align=center, inner sep=3pt}]
  \node[obs] (L) {Lei};
  \node[obs, right=of L] (A) {Aplicação\\aérea};
  \node[lat, right=of A] (X) {Exposição\\ambiental};
  \node[obs, right=of X] (Y) {Desfecho\\ao nascer};
  \node[obs, below=of A] (T) {Aplicação\\terrestre};
  \node[lat, below=of X] (O) {Exposição\\ocupacional};
  \node[obs, above=of A] (P) {Produção\\e emprego};
  \node[obs, above=of P] (C) {Composição\\das mães};
  \node[obs, below=of Y] (F) {Perda\\fetal};
  \node[obs, right=of Y] (G) {Choque\\estadual};
  \draw[->, very thick] (L) -- (A);
  \draw[->, very thick] (A) -- (X);
  \draw[->, very thick] (X) -- (Y);
  \draw[->] (L) |- (T);
  \draw[->] (T) -- (X);
  \draw[->] (T) -- (O);
  \draw[->] (O) -- (Y);
  \draw[->] ([xshift=1.5mm]L.north) |- (P);
  \draw[->] ([xshift=-1.5mm]L.north) |- (C);
  \draw[->] (P) -| ([xshift=-2mm]Y.north);
  \draw[->] (C) -| ([xshift=2mm]Y.north);
  \draw[->] (X) -- (F);
  \draw[->, dashed] (G) -- (Y);
\end{tikzpicture}
\caption{Grafo causal resumido. Em traço grosso, o caminho-alvo. Caixas
tracejadas: variáveis não observadas. Seta tracejada: ameaça à identificação.}
\label{fig:dag}
\end{figure}

\begin{enumerate}
\item \textbf{O caminho-alvo:} lei, menos avião, menos deriva, melhor
nascimento.
\item \textbf{Substituição ambiental:} a aplicação terrestre também expõe os
vizinhos, em menor grau. O efeito estimado é líquido dela.
\item \textbf{Substituição ocupacional:} mais trabalhadores aplicam, e pais
expostos podem afetar a gestação.
\item \textbf{Renda:} a lei altera custo, produção e emprego, que alteram a
renda das famílias.
\item \textbf{Composição:} a lei pode mudar quem mora ou engravida nos
municípios atingidos.
\item \textbf{Seleção para nascer vivo:} se a perda fetal cai, nascem fetos
mais frágeis, e o peso médio dos nascidos vivos fica atenuado.
\end{enumerate}

Para atribuir o efeito ao primeiro caminho, a avaliação mostra, nesta ordem: que
o avião de fato deixou de voar (o primeiro estágio); que produção, emprego e
perfil das mães não se moveram; e o que aconteceu com as intoxicações de
trabalhadores e com a perda fetal. É a análise de mecanismo que a teoria da
mudança torna possível.
